\section{Introduction}
Material values provide a compact vocabulary for discussing chess decisions. They make it possible to summarize a trade, describe a material imbalance, or compare two configurations without displaying an entire search tree. A familiar numerical table is nevertheless a substantial compression: it assigns one number to a piece type while the available moves depend on location, pawn structure, king safety, and the remaining material. The present study examines a restricted empirical question: what relative material coefficients are needed to explain an engine's observed move choices when a transparent reward model is fitted separately or conditionally across different populations of positions?

This question concerns an observable behavioral regularity, rather than a reconstruction of every component of the engine. An engine returns a selected move after search. A model fitted to those selections can describe which immediate action features are associated with the choices, but its coefficients need not equal the engine's internal evaluation parameters. The distinction matters especially when a move offers material temporarily, anticipates a recapture, or exchanges a positional advantage for a future tactical opportunity. A coefficient model that sees only the resulting one-ply feature difference cannot represent that continuation. Its numerical weights are therefore conditional on both its feature vocabulary and its deliberately limited decision horizon.

\subsection{From observed choices to relative values}
Inverse reinforcement learning (IRL) frames behavior as evidence about a reward function \cite{ng2000,ziebart}. In a full sequential formulation, evaluating a proposed reward also involves the consequences of a policy over time. Here we specialize to a horizon of one decision and enumerate every legal move from each observed state. The resulting maximum-entropy distribution is a conditional logit model whose score is linear in action-feature differences. This formulation is small enough that its fitted coefficients, objective, candidate sets, and uncertainty calculations can all be inspected directly.

We define an ``IRL-estimated relative piece value'' as a fitted material reward coefficient divided by the fitted pawn coefficient under a specified model and reference population. This definition has two advantages for the intended analysis. First, each reported value has an explicit statistical meaning. Second, the same convention can be applied to different subsets without silently replacing the target with an engine evaluation score or a predicted game outcome. It also imposes a limitation: a ratio below one describes a fitted immediate-choice coefficient, not a recommendation to exchange that piece for a pawn. Ratios additionally depend on the denominator, so we retain unnormalized coefficients when interpreting changes.

\subsection{Two complementary descriptions of position context}
We first summarize opening, middlegame, and endgame positions. These labels are familiar descriptions of a game's progression, but they require operational definitions for reproducible analysis. Our endgame rule uses weighted remaining non-pawn material; opening membership uses move number after excluding endgames; all remaining positions are assigned to the middlegame. We report the exact cutoffs rather than assuming that the labels have one universally accepted meaning. A pooled contextual model uses continuous material phase and file openness, and its fitted coefficients are averaged at stage-specific reference contexts. Separate stage fits provide a comparison with less cross-stage information sharing.

We then ask a distinct question about how many pieces remain. A board with few occupied squares need not have the same material composition as another board with the same count. Moreover, weighted non-pawn material can fall while many pawns remain. We therefore add independent fits for sparse, intermediate, and dense positions, counting all non-king pieces of both colors, including pawns. This extension tests whether the descriptive pattern seen across stages also appears under a simpler occupancy-based partition. It is not a controlled intervention that removes pieces from an otherwise identical position.

The two partitions intentionally overlap. Reporting that overlap is necessary to understand what the second analysis contributes. If nearly all sparse positions are already endgames, agreement between the two tables is partly shared evidence. Conversely, a mixed intermediate group can reveal that a stage label and a piece-count threshold select different combinations of tactical opportunities. We report group sizes, stage composition, and the number of states whose legal alternatives actually distinguish each material feature. These quantities are directly relevant to interpreting the value table, especially when a visually large coefficient is supported by only a few informative choices.

\subsection{Research questions and contribution}
The investigation has four empirical questions. First, how do pawn-normalized coefficients differ across the three stage reference populations of a pooled contextual model? Second, how much do the estimates change when each stage or piece-count group is fitted independently? Third, which differences remain visible under the tested choices of regularization and group boundary? Fourth, how do the estimates respond to Stockfish strength settings 1500, 2000, and 2500 when the observed boards are held fixed? These questions require consistent definitions, actual fits, uncertainty estimates, and reporting of alternatives that do not necessarily agree; they do not require claiming a new IRL algorithm.

The dataset contains 300 Stockfish 18 self-play games with fixed engine settings and controlled early exploration. After excluding forced decisions and exact repeated board states according to a fixed split precedence, 18,351 positions remain. Game-level train, validation, and test partitions keep later observations from a game out of another partition. Hyperparameters are selected on validation data. The principal outcomes are the learned relative coefficients and their dependence on the modeling choices. The original test partition is preserved and is not used for coefficient estimation or hyperparameter selection.

Our contributions are an explicit horizon-one estimand for stage-specific relative material values; pooled and independently fitted estimates with game-bootstrap intervals; a new comparison across remaining-piece-count groups using paired game resampling; and an auditable account of support, normalization, and sensitivity. The additional group analyses reuse the original collection and were motivated after inspecting the original stage results. They are therefore exploratory extensions, not preregistered confirmatory tests or independent replications. Reporting this sequence is part of making the strength of the evidence clear.

\subsection{Scope of the claims}
Our conclusions apply to one engine version, search budget, state distribution, and feature representation. Stronger search, human demonstrations, or a reply-aware reward model may yield different coefficients. A statistically separated ratio cannot establish a causal effect of game progression or an intrinsic material price. We use the term relative value only with this model-conditional meaning. The remaining sections locate the study among chess evaluation, behavioral imitation, and IRL; define the estimator and data; present stage and piece-count results; and document their limits and reproduction.
